(1) 静力平衡给 $dp/dz=-\rho g=-mng$，理想气体满足 $p=nk_BT$。因 $T$ 恒定，得到 $d\ln n/dz=-mg/(k_BT)$，积分后 $n(z)=n_0e^{-z/H}$、$p(z)=p_0e^{-z/H}$，其中 $p_0=n_0k_BT$、$H=k_BT/(mg)$，$H$ 的量纲为长度。
(2) 整层大气的单位面积分子总数 $N/A=\int_0^\infty n(z)dz=n_0H$，柱质量 $M/A=mn_0H=n_0k_BT/g=p_0/g$。因而底部压强恰等于其上方整层大气的重量除以面积，既验证微分解也表明假设上边界取无穷远的情况下总柱质量有限。
(3) 任取一分子的高度概率密度是 $f(z)=n(z)/(n_0H)=e^{-z/H}/H$，$z\ge0$。通过 $\int_0^\infty z e^{-z/H}dz=H^2$ 得 $\langle z\rangle=H$，于是 $\langle U_g\rangle=mg\langle z\rangle=mgH=k_BT$。结果与 $m$ 无关是因为较重粒子的尺度高度相应更小。
(4) 高于 $z_*$ 的比例为 $P(z>z_*)=\int_{z_*}^\infty f(z)dz=e^{-z_*/H}$，所以一半分子位于其上的中位高度满足 $z_{1/2}=H\ln2$。真实大气的温度变化、可变重力、组分分离、对流、湍流及有限外逸会使该指数廓线失效；它是理想静态等温模型而不是整层地球大气的普适精确描述。
